\section{Affine labels with a prescribed tolerance}
\label{sec:affine}

We construct a finite minimum of affine functions whose minimizing slopes
approximate the exponents of a positive point.  The permitted error is
prescribed separately at each slope.  Allowing this prescription to be
discontinuous will let us impose sign conditions on finitely many linear
forms in the application.

For an integer $m\geq 1$, $u\in\R^m$, and $0<h<1$, write
\[
  X_h(u)=(h^{u_1},\ldots,h^{u_m}).
\]
A \emph{label} is a pair $(r,c)$ consisting of a vector $r\in\R^m$,
independent of $h$, and a real-valued offset $c(h)$ defined for all
sufficiently small positive $h$.  For a finite nonempty family $\Lambda$
of labels, write its members as $\lambda=(r_\lambda,c_\lambda)$.
For fixed sufficiently small $h>0$ and $x\in\R^m$, a label
$\lambda\in\Lambda$ is \emph{active at $x$} if
\[
  r_\lambda\cdot x+c_\lambda(h)
  =\min_{\eta\in\Lambda}\bigl(r_\eta\cdot x+c_\eta(h)\bigr).
\]
Thus every minimizing label is active when a tie occurs.  Whenever finitely
many families are combined, we restrict their offsets to a common interval
$0<h<h_*$ on which they are all defined.

\begin{lemma}[Affine approximation with arbitrary tolerance]
\label{lem:affine}
Let $m\geq 1$ be an integer, let $a<b$ be real numbers, and let
$E:[a,b]^m\to(0,\infty)$ be any function.  There is a finite nonempty
family $\Lambda$ of labels $\lambda=(r_\lambda,c_\lambda)$, with
$r_\lambda\in[a,b]^m$ and $c_\lambda(h)=o(h^a)$ as $h\to0+$, such
that the following holds.
For every sequence $h_\nu\to0+$ and every sequence $p^\nu\in\R^m$
converging to $p\in[a,b]^m$, if a fixed label $\lambda\in\Lambda$ is
active at $X_{h_\nu}(p^\nu)$ for every $\nu$, then
\[
  \norm{p-r_\lambda}_\infty<E(r_\lambda).
\]
No regularity of $E$ is assumed, and the approximating vectors $p^\nu$
need not belong to $[a,b]^m$.
\end{lemma}

Two aspects of the statement will be used separately. The slopes form a
fixed finite set, rather than a set that changes with $h$. Also, the
conclusion concerns every label that stays active along a sequence,
including labels tied with other minimizers. The freedom for $p^\nu$ to
leave the cube is needed by the induction: a projected sequence may
approach a face of a smaller cube from either side.

\begin{proof}
We induct on $m$, using only positive dimensions smaller than $m$.
The largest coordinate of $X_h(p)$ is $h^{\min_i p_i}$, so we first
construct a family for each fixed value of the smallest exponent.
Within that family, negative offsets compensate for raising slope entries
on the smaller coordinates of $X_h(p)$.
Compactness extends the family's guarantee to nearby minimum
values, and a finite interval construction then selects among these
families. When $m=1$, the first step
uses singleton families and invokes no induction hypothesis; the remaining
steps therefore also establish the base case.

\medskip
\noindent\emph{Step 1: construction at a fixed minimum.}
Set $\mu(u)=\min_i u_i$, write $\mathbf 1=(1,\ldots,1)$, and fix
$t\in[a,b]$.  We will construct a finite
family $\Lambda_t$ whose slopes all satisfy $\mu(r)=t$, containing
$(t\mathbf 1,0)$, with the following property: the conclusion of the lemma
holds for activity within $\Lambda_t$ whenever the limiting vector
$p\in[a,b]^m$ satisfies $\mu(p)=t$.

If $m=1$ or $t=b$, take $\Lambda_t=\{(t\mathbf 1,0)\}$.  In either
case, a limit in the cube with minimum $t$ must equal $t\mathbf 1$, so
the asserted property holds.  Assume henceforth in this construction that
$m>1$ and $t<b$.

A label has \emph{type $k$} if precisely $k$ entries of its slope are
strictly greater than $t$, with all other entries equal to $t$.  Type $0$
consists of $(t\mathbf 1,0)$.  Construct types $k=1,\ldots,m-1$ in
increasing order.  After type $k-1$ has been constructed, choose $\beta_k$
so that
\begin{gather*}
  t<\beta_k<b,\qquad
  \beta_k<\beta_{k-1}\quad\text{if }k>1,\\
  \beta_k-t<E(r)
  \quad\text{for every slope $r$ of type $k-1$.}
\end{gather*}
The already constructed type has finitely many slopes, and its tolerance
values are all positive, so such a choice is possible. Only those finitely
many values are compared; we do not take a positive lower bound for $E$
on an entire cube.

For each $k$-element subset $U\subset\{1,\ldots,m\}$ and each
$s\in[\beta_k,b]^U$, let $r(s)\in[a,b]^m$ have entries $s$ on $U$ and
$t$ off $U$.  Apply the induction hypothesis in dimension $k$ to
$[\beta_k,b]^U$, with tolerance $s\mapsto E(r(s))$; the coordinates of
$U$ are taken in their natural order.  For every resulting label $(s,d)$,
where $d(h)=o(h^{\beta_k})$, put the label
\[
  r=r(s),\qquad c(h)=-h^{\beta_k}+d(h)
\]
in type $k$.  Taking all such coordinate sets and types gives a finite
nonempty family $\Lambda_t$.  Each of its offsets is $o(h^t)$.
The negative term $-h^{\beta_k}$ favors raising slope entries when
the corresponding coordinates of $x$ are small relative to $h^t$; the decreasing
$\beta_k$ make this preference stronger at each successive type.

\medskip
\noindent\emph{Step 2: verification at the fixed minimum.}
For the construction with $m>1$ and $t<b$, suppose $h_\nu\to0+$,
$p^\nu\to p\in[a,b]^m$, and $\mu(p)=t$.
Fix a label $(r,c)$ of type $k$ that is active within $\Lambda_t$ along this
sequence, and let $U$ be its set of raised coordinates.  All asymptotic
comparisons in this verification are along the sequence, and we write
$h=h_\nu$. Whenever a limiting exponent is strictly above or below a
threshold, convergence leaves a fixed positive exponent gap for all large
$\nu$. We use these gaps, not a replacement of $h^{p_i^\nu}$ by
$h^{p_i}$: convergence of the exponents alone would not justify that
replacement.

If $k>0$, comparison with $(t\mathbf 1,0)$ gives
\[
  \sum_{i\in U}(r_i-t)h^{p_i^\nu}
  \leq h^{\beta_k}-d(h)
  =(1+o(1))h^{\beta_k}.
\]
Every coefficient on the left is at least $\beta_k-t>0$.  If some
$p_i<\beta_k$ for $i\in U$, then, for a fixed $\delta>0$ and all
sufficiently large $\nu$, $p_i^\nu\leq\beta_k-\delta$; dividing the
display by $h^{\beta_k}$ would yield a term tending to infinity.
Consequently
\[
  p_i\geq\beta_k\quad(i\in U).
\]
Among labels constructed from this same $U$, the contributions off $U$
and the offset $-h^{\beta_k}$ are common.  Thus $(s,d)$ is active in the
family supplied by induction.  Its limiting exponent vector belongs to
$[\beta_k,b]^U$, so the induction hypothesis yields
\[
  \max_{i\in U}\abs{p_i-r_i}<E(r).
\]
This is precisely where unrestricted approximating vectors are needed:
$p_U^\nu$ may lie outside $[\beta_k,b]^U$ even though its limit lies
inside that cube. In particular, no strict inequality is required when a
limiting coordinate equals $\beta_k$ or $b$.

We next control the coordinates off $U$, including type $0$ by taking
$U=\varnothing$.  Suppose $k<m-1$.  We claim that
\[
  p_i\leq\beta_{k+1}\quad(i\notin U).
\]
If instead $p_j>\beta_{k+1}$ for some $j\notin U$, then every limiting
exponent on $U\cup\{j\}$ is strictly greater than $\beta_{k+1}$.
For $k>0$, this uses $p_i\geq\beta_k>\beta_{k+1}$ on $U$; for $k=0$
there are no such coordinates to check.  Choose any type $k+1$ label
whose raised set is $U\cup\{j\}$.  Its affine value minus the value of
$(t\mathbf 1,0)$ is
\[
  -h^{\beta_{k+1}}+o(h^{\beta_{k+1}}),
\]
because $h^{p_i^\nu}=o(h^{\beta_{k+1}})$ on every raised coordinate.
The corresponding difference for the supposed active label is zero if
$k=0$, and is bounded below by
$-h^{\beta_k}+o(h^{\beta_k})$ if $k>0$.
Since $h^{\beta_k}=o(h^{\beta_{k+1}})$ in the latter case, the chosen
competitor has a strictly smaller value for all sufficiently large $\nu$,
a contradiction.  This proves the claim.  Because $p_i\geq t$ and
$r_i=t$ off $U$, it follows that
\[
  \abs{p_i-r_i}\leq\beta_{k+1}-t<E(r)
  \quad(i\notin U),
\]
where the strict inequality was imposed when constructing type $k+1$.

If instead $k=m-1$, the sole coordinate off $U$ has limiting exponent
exactly $t$: all coordinates on $U$ have exponent at least
$\beta_k>t$, and $\mu(p)=t$.  Its error is therefore zero.  Combining
the bounds proves $\norm{p-r}_\infty<E(r)$ in all cases, including
types $0$ and $m-1$.  No strictness was required at a threshold
$p_i=\beta_k$.

\medskip
\noindent\emph{Step 3: extension to nearby minima.}
Keep $t$ and its finite family $\Lambda_t$ fixed.  For a label
$\lambda=(r,c)$ in that family, let $A_\lambda\subset[a,b]^m$ be the
set of all limits of sequences $p^\nu\in\R^m$ for which
$h_\nu\to0+$ and $\lambda$ is active within $\Lambda_t$ at
$X_{h_\nu}(p^\nu)$.
This set is closed.  Indeed, if $z_n\in A_\lambda$ and $z_n\to z$,
choose from a realizing sequence for $z_n$ a point $u_n$ and parameter
$h_n$ such that
\[
  0<h_n<1/n,\qquad \norm{u_n-z_n}_\infty<1/n,
\]
and $\lambda$ is active at $X_{h_n}(u_n)$.  Then $u_n\to z$, so
$z\in A_\lambda$. This diagonal selection uses actual active points;
it does not pass activity through a continuity assertion about $c(h)$.

The set
\[
  B_t=\bigcup_{\lambda=(r,c)\in\Lambda_t}
  \bigl\{p\in A_\lambda:\norm{p-r}_\infty\geq E(r)\bigr\}
\]
is therefore compact.  Here each $E(r)$ is a fixed real number;
continuity of $E$ is unnecessary.  The fixed-minimum verification shows
that $B_t$ contains no point with $\mu(p)=t$.  Since $\mu$ is continuous,
there is $\rho_t>0$ such that $B_t$ contains no point with
\[
  \mu(p)\in I_t:=(t-\rho_t,t+\rho_t).
\]
For example, if $B_t$ is nonempty, take $\rho_t$ smaller than the
positive minimum of $\abs{\mu(p)-t}$ on $B_t$; if it is empty, any
positive radius initially suffices.  Thus the limiting active-label
conclusion holds whenever $p\in[a,b]^m$ and $\mu(p)\in I_t$.

Shrink $\rho_t$, if necessary, so that every local offset also satisfies
\begin{equation}
  c(h)=o(h^{t+\rho_t}).
  \label{eq:affine-local-decay}
\end{equation}
For a nonsingleton family it suffices to require
$t+\rho_t<\min_{1\leq k\leq m-1}\beta_k$; for a singleton family the
offset is zero. The stronger decay uses the constructed form
$c(h)=-h^{\beta_k}+o(h^{\beta_k})$ with $\beta_k>t$; it would not follow
from $c(h)=o(h^t)$ alone. Shrinking the radius preserves the guarantee
just proved.
We have now obtained the desired approximation near every possible value
of $\mu(p)$.  It remains to make a finite combined family select an
appropriate one of these neighborhoods.

\medskip
\noindent\emph{Step 4: finite interval construction and combination.}
The local families already approximate correctly near their own minimum
levels. We now give each family a common offset shift so that a globally
active label must come from a suitable neighborhood. The overlap between
neighboring neighborhoods is important: it permits equality at the
transition cuts.

Choose a finite subcover of $[a,b]$ from the intervals $I_t$.  Remove
intervals contained in another selected interval, retaining one copy if
two coincide.  Order the remaining centers as
\mbox{$a\leq t_1<\cdots<t_N\leq b$}, and write
$I_{t_j}=(L_j,R_j)$.  Noncontainment forces the orders of the left and
right endpoints to agree; since $t_j=(L_j+R_j)/2$, both orders agree
with that of the centers.  Hence
\[
  L_1<\cdots<L_N,\qquad R_1<\cdots<R_N.
\]
Successive intervals overlap.  Otherwise $R_j\leq L_{j+1}$, and a
point of $[R_j,L_{j+1}]\subset(t_j,t_{j+1})$ would be covered by none
of the selected intervals.  The first interval contains $a$ and the
last contains $b$, by the endpoint orders, coverage, and the fact that
all centers lie in $[a,b]$.

For $1\leq j<N$, choose
\[
  q_j\in
  \bigl(\max\{t_j,L_{j+1}\},\,
        \min\{t_{j+1},R_j\}\bigr).
\]
This open interval is nonempty: the two intervals overlap, each contains
its own center, and $t_j<t_{j+1}$.  Thus $q_j$ belongs to both adjacent
intervals and lies strictly between their centers.  Setting $q_0=a$ and
$q_N=b$, we obtain
\begin{equation}
  t_j<q_j<t_{j+1}\quad(1\leq j<N),\qquad
  [q_{j-1},q_j]\subset I_{t_j}\quad(1\leq j\leq N).
  \label{eq:affine-cuts}
\end{equation}
If $N=1$, no cuts are needed and the sole selected interval contains
the whole of $[a,b]$.

Take the union of the selected families, replacing every offset $c(h)$
at level $t_j$ by
\begin{equation}
  \widetilde c(h)=c(h)-\sum_{\ell=1}^{j-1}h^{q_\ell}.
  \label{eq:affine-shifts}
\end{equation}
The empty sum at $j=1$ is zero.  Every new offset is $o(h^a)$:
Equation~\eqref{eq:affine-local-decay} applies with $t_j+\rho_{t_j}>a$,
and every shift exponent satisfies $q_\ell>t_\ell\geq a$.
Within a local family this common shift does not change activity. Between
adjacent families the new term $-h^{q_j}$ selects the family on the
appropriate side of the cut. Figure~\ref{fig:interval-gluing} shows the
interval cover and the closed cut cells assigned to the selected families.

\begin{figure}[!bp]
\centering
\begin{tikzpicture}[x=0.88cm,y=0.75cm,font=\small]
\draw[-{Stealth}] (-1.3,0)--(11.5,0) node[right] {$\mu(p)$};
\draw[densely dotted,gray] (3.8,0.2)--(3.8,3.45);
\draw[densely dotted,gray] (6.2,0.2)--(6.2,3.45);
\draw[blue!65!black,line width=1pt] (-1,3)--(5,3);
\draw[blue!65!black,line width=4pt] (0,3)--(3.8,3);
\draw[teal!75!black,line width=1pt] (2.4,2)--(7.6,2);
\draw[teal!75!black,line width=4pt] (3.8,2)--(6.2,2);
\draw[violet!75!black,line width=1pt] (5,1)--(11,1);
\draw[violet!75!black,line width=4pt] (6.2,1)--(10,1);
\draw[blue!65!black,fill=white] (-1,3) circle (2pt);
\draw[blue!65!black,fill=white] (5,3) circle (2pt);
\draw[teal!75!black,fill=white] (2.4,2) circle (2pt);
\draw[teal!75!black,fill=white] (7.6,2) circle (2pt);
\draw[violet!75!black,fill=white] (5,1) circle (2pt);
\draw[violet!75!black,fill=white] (11,1) circle (2pt);
\fill[blue!65!black] (2,3) circle (2pt);
\fill[teal!75!black] (5,2) circle (2pt);
\fill[violet!75!black] (8,1) circle (2pt);
\node[above] at (2,3.12) {$I_{t_1}$};
\node[above] at (5,2.12) {$I_{t_2}$};
\node[above] at (8,1.12) {$I_{t_3}$};
\draw (0,0.08)--(0,-0.08) node[below] {$a=q_0$};
\draw (2,0.08)--(2,-0.08) node[below] {$t_1$};
\draw (3.8,0.08)--(3.8,-0.08) node[below] {$q_1$};
\draw (5,0.08)--(5,-0.08) node[below] {$t_2$};
\draw (6.2,0.08)--(6.2,-0.08) node[below] {$q_2$};
\draw (8,0.08)--(8,-0.08) node[below] {$t_3$};
\draw (10,0.08)--(10,-0.08) node[below] {$b=q_3$};
\end{tikzpicture}
\caption{Finite gluing, schematically with three selected families.
The open intervals $I_{t_j}$ are the neighborhoods where the local
families approximate correctly. Each thick closed segment
$[q_{j-1},q_j]$ lies strictly inside its corresponding open interval.
For a label from family $j$ that stays active along a convergent sequence,
the common shifts in \eqref{eq:affine-shifts} force the limiting exponent
minimum $\mu(p)$ into that closed segment. Either adjacent family remains
valid when $\mu(p)$ equals a cut.}
\label{fig:interval-gluing}
\end{figure}


Suppose a fixed label from level $t_j$ is active in this combined family
along $h_\nu\to0+$ and $p^\nu\to p\in[a,b]^m$.  Write $c(h)$ for
its original local offset and again abbreviate $h=h_\nu$.  We claim
\[
  q_{j-1}\leq\mu(p)\leq q_j.
\]
If $j<N$ and $\mu(p)>q_j$, every $h^{p_i^\nu}$ is $o(h^{q_j})$.
The local offset is also $o(h^{q_j})$, by
Equation~\eqref{eq:affine-local-decay} and $q_j\in I_{t_j}$.
Subtracting the value of the label
$(t_{j+1}\mathbf 1,0)$ with its shift, the value difference is
\[
  (r-t_{j+1}\mathbf 1)\cdot X_h(p^\nu)+c(h)+h^{q_j}
  =h^{q_j}+o(h^{q_j})>0
\]
for all sufficiently large $\nu$, contradicting activity.

If $j>1$ and $\mu(p)<q_{j-1}$, subtract instead the shifted
label $(t_{j-1}\mathbf 1,0)$.  The difference is
\[
  (r-t_{j-1}\mathbf 1)\cdot X_h(p^\nu)+c(h)-h^{q_{j-1}}.
\]
Every coefficient in the scalar product is at least
$t_j-t_{j-1}>0$.  For some coordinate $i$, $p_i<q_{j-1}$, so
$h^{p_i^\nu}/h^{q_{j-1}}\to\infty$.  On the other hand,
$c(h)=o(h^{q_{j-1}})$ because $q_{j-1}<t_j<t_j+\rho_{t_j}$.
The displayed difference divided by $h^{q_{j-1}}$ therefore tends to
$+\infty$, again contradicting activity.  The outer bounds when $j=1$
or $j=N$ follow from $\mu(p)\in[a,b]$.

By Equation~\eqref{eq:affine-cuts}, the claimed bounds imply
$\mu(p)\in I_{t_j}$, including equality at either cut.  A label active
in the combined family is active within its original local family,
since all labels at that level received the same shift.  The guarantee
for $I_{t_j}$ now gives $\norm{p-r}_\infty<E(r)$.
All contradictions above produced a strictly smaller competing value,
so the argument applies to every active label even at ties.  This
completes the induction.
\end{proof}

The sequential conclusion gives a uniform estimate on the entire cube,
with the same fixed finite family. This is the form used for reaction
networks.

\begin{corollary}[Uniform activity bound]\label{cor:uniform-affine}
Let $m\ge1$, $a<b$, and $E:[a,b]^m\to(0,\infty)$ be as in
Lemma~\ref{lem:affine}. Fix a finite nonempty family
$\{(r_\lambda,c_\lambda):\lambda\in\Lambda\}$ supplied by that Lemma.
There is $h_0>0$ such that, for every $0<h<h_0$ and every
$p\in[a,b]^m$, each label active at $X_h(p)$ satisfies
\[
 \norm{p-r_\lambda}_\infty<E(r_\lambda).
\]
The family is fixed before $h_0$ is chosen.
\end{corollary}

\begin{proof}
Otherwise there are $h_\nu\to0+$, points $p^\nu\in[a,b]^m$, and
labels $\lambda_\nu$ active at $X_{h_\nu}(p^\nu)$ for which
\[
 \norm{p^\nu-r_{\lambda_\nu}}_\infty
 \ge E(r_{\lambda_\nu}).
\]
By finiteness of $\Lambda$ and compactness of $[a,b]^m$, a subsequence
has a fixed label $\lambda_\nu=\lambda$ and $p^\nu\to p\in[a,b]^m$.
Passing to the limit in the displayed inequality contradicts the strict
inequality in Lemma~\ref{lem:affine}. The value $E(r_\lambda)$ is fixed
on this subsequence, so no regularity of $E$ is used.
\end{proof}

In Section~\ref{sec:dynamics}, we choose the tolerance so that every
nonzero dot product of an active slope with a prescribed vector forces
the exponent vector to have a dot product of the same sign. The freedom
to choose a discontinuous $E$ permits this requirement near the corresponding
hyperplanes.
